\markboth{}{Preface}
\cleardoublepage
\phantomsection
\addcontentsline{toc}{chapter}{Preface}
\chapter*{Preface}

\textsc{Pāli\,Platform\,1} had been released in January 2020. It might be only I who used it substantially to produce Pāli learning books.\footnote{\emph{Pāli for New Learners} and \emph{Pāli Text Reading:\,A Handbook}. See \url{bhaddacak.github.io/pnl} and \url{bhaddacak.github.io/ptr} respectively.} After I finished the books, I had spent several months (9 or so) to rewrite \textsc{Pāli\,Platform}, to make it modernized in look---more user-friendly, to add more features, and to fix many bugs.\footnote{I moved from Java 8 to Java 11 and replaced Swing UI with JavaFX. I thought at first I would use Kotlin, but it is still more difficult to work with. Java is still the best friend I have. It is not the language itself is great, but the whole Java ecosystem makes programming with Java enjoyable: enormous available libraries, fast compile time, excellent API documentation, flexible and uncomplicated deployment scheme, etc. Moreover, as I have worked with JavaFX so far, this GUI is the best among all cross-platform GUI toolkits. It can run and looks the same in all platforms without any adjustment. Some might say Java is slow, at least in comparison with C and C++. As you shall see by using the program, it just needs start-up time (big programs written in C or C++ also need start-up time, even longer than typical Java applications). After that, the difference is barely noticeable. Many programmers dislike the verbosity of Java syntax, but I see this as its strength. Verbosity makes us more mindful/careful and helps us make less mistakes. I do not mind typing more when coding, because pounding the keyboard is a kind of fun.} Around the end of 2022, \textsc{Pāli\,Platform\,2} had been released, bundled with the first printable manual.\footnote{Known as \emph{Pāli Platform 2:\,The Official Manual}}

Then at the start of 2025 or shortly before that, we had seen the launch of \textsc{Pāli\,Platform\,3} with several new features. So, a substantial revision of the old manual was needed.

Recently, \textsc{Pāli\,Platform\,4} has been launched at the start of 2026. This version has several improvements and additions, notably with a number of useful Sanskrit tools. Now we can say that \textsc{Pāli\,Platform} is the most comprehensive program for learning Pāli and Sanskrit. The cost of this is the inevitable complexity. So, after the new release has run smoothly, I pick up the manual and make it up-to-date.\footnote{Many of screenshots are newly updated, but some which their main area is unaffected are retained. There can be a small difference in the tool bar.}

After a stable release of \textsc{Pāli\,Platform\,4}, all older versions\footnote{\textsc{Pāli\,Platform\,1} had been removed from the public since the launch of \textsc{Pāli\,Platform\,2}. After that, \textsc{Pāli\,Platform\,2}, by its new project code `\textsc{Pāli\,Platform\,Classic}', had been removed from the public as well. The last release of \textsc{Pāli\,Platform\,3} is now marked as deprecated and will be deleted soon.} will be discarded to reduce the maintaining cost.

The manual is divided into eight parts. The first part, Essential Starter, is crucial. It is supposed to be read deliberately. The main purpose of this is to help the users start the program successfully in various contexts, and to provide an initial guidance and troubleshooting.

The second part is about dictionaries. Since \textsc{Pāli\,Platform\,3}, Digital Pāḷi Dictionary (DPD) has played a major role in word look-up. So, this part have a chapter on the \hyperref[chap:dpd]{DPD}. Recently, explanations on Sanskrit dictionaries is also added in Chapter \ref{chap:dict} (other Sanskrit tools can be seen in their part below).

The third part is all about Pāli grammatical tools. It is enough to just go through the part quickly and come back when necessary. Several chapters are short. The longest one is Chapter \ref{chap:prosody} on Prosody, which needs an elaborate treatment.

The fourth part is on Pāli collections. You will learn how to find a document in the corpora we have, and how to use text viewers. The content of this part is fundamental to the whole program. So, it should be read carefully.

The fifth part is about advanced search tools. It has two sizable chapters on Lucene Finder and Term Lister. These two chapters have to be looked closely because of the complexity of the tools. The functionality of this part has determined the existence of the program from the start.

The sixth part is newly added. This consists of a number of Sanskrit tools (apart from dictionaries). To Pāli learners, these tools extend the horizon of the knowledge considerably. And to Sanskrit learners, they are really handy and helpful.

The seventh part talks about the rest of the all above. You will learn various tools for various purposes, such as the program's Text Editor, and Sentence Reader. This Sentence Reader, together with its companion tool, Sentence Manager, is an innovative tool that can help the learners read Pāli texts more conveniently. Furthermore, translations can be added to the texts at sentence level by this toolset. Tools for technical users is also explained to this part. This can be difficult for some, but it may be a piece of cake for power users. The last chapter is a short treatment of regular expression. For the topic itself is big and beyond our main concern, what I can do is just a survival introduction.

The last part is appendices containing three chapters. The first one is about the conversion between Sanskrit and Pāli words. This can help Pāli learners enhance their ability of finding meaning of words parallel in both languages. The second appendix is the list of prosodic patterns, not present in Chapter \ref{chap:prosody} (Prosody) for its complexity. It is put here for advanced learners who may ask for this. The last one is GNU Free Documentation License. This means the manual is fully opened to the public (with conditions stated in the license).

My target readers of this manual are those who want to make use of \textsc{Pāli\,Platform} to its full capacity, both for learning and researching purpose. Some basic knowledge of Pāli (and Sanskrit) is helpful, particularly the terminology used in the field. For the fundamental of the language, see the books mentioned. Users outside the sphere of computational technology may skip computer-related technical terms that appear unexplained, or else they can surely find an explanation in Wikipedia or AI chat/search engines.

This manual is intended to be bundled with \textsc{Pāli\,Platform\,4}, base version 4.2.0. If errors, lurking somewhere in the manual, books, and program, are found, please kindly report them to the author.\footnote{\texttt{bhaddacak} at \texttt{proton} dot \texttt{me}}

